\chapter{分析：把直觉变成可靠的定理}

\section{山路上的一道怪题}

暑假里，你和家人去爬山。山脚的海拔是 $100$ 米，山顶的海拔是 $1500$ 米。你们早上八点出发，走走停停，中午十二点半到达山顶。途中你们休息过，折返过一段路去找掉在地上的帽子，速度时快时慢，总之走得毫无规律。

回家后，有人问你一个奇怪的问题：“在爬山过程中，你们一定有一个瞬间，海拔恰好是 $618$ 米。你敢保证吗？”

你凭直觉会说：“当然！从 $100$ 米到 $1500$ 米，一路上海拔是连续变化的，总不能从 $617$ 米‘跳’到 $619$ 米吧？”

这个直觉听起来天衣无缝。但请先别急着点头，再听两个问题：

\begin{itemize}
  \item 登山记录显示，早上九点你的海拔是 $500$ 米，上午十一点是 $800$ 米。那么是否一定存在某个时刻，海拔恰好是 $600.001$ 米？
  \item 一支温度计从清晨到正午，温度从 $10^\circ$C 升到 $25^\circ$C。是否一定有某个时刻，温度恰好等于 $\sqrt{500}$ 摄氏度？注意，$\sqrt{500}$ 是一个无限不循环小数，任何温度计的刻度都读不出它。
\end{itemize}

第一个问题你可能仍然答“是”，但底气已经开始动摇：海拔真的可以无限精确地谈吗？第二个问题更麻烦——一个在任何温度计上都显示不出来的数，凭什么说它“被恰好经过”过？

这一章的任务，就是把“连续变化不会跳过任何中间值”这句模糊的感觉，锻造成一条可以严格证明、可以放心使用的定理。这门手艺叫\textbf{分析}：它的目标不是算得更快，而是把直觉变成可靠的定理。数学家们为此折腾了整整两百年——他们的直觉被骗过不止一次，骗局还非常精彩。

\section{直觉第一次被骗}

先给“连续”一个最朴素的理解：函数的图像是一条“不抬笔就能画完”的曲线。这个理解有用，但它靠不住。历史上它至少在三件事上栽过跟头。

第一件事：你以为“不跳过中间值”是显然的，那请看这个函数：
\[
f(x)=\begin{cases}1, & x\ge 0,\\ -1, & x<0.\end{cases}
\]
它在 $x=-1$ 处取值 $-1$，在 $x=1$ 处取值 $1$，中间值 $0$ 却永远取不到——图像在 $x=0$ 处“跳了台阶”。所以“不跳过中间值”这件事，真的和“不抬笔”有关，不是白送的。

第二件事更隐蔽。我们知道 $1^2=1<2$，$2^2=4>2$，所以 $x^2=2$ 的解 $\sqrt{2}$ 应该“夹在 $1$ 和 $2$ 之间被经过”。可是请你回忆前面的章节：$\sqrt{2}$ 不是有理数。假如世界上只有有理数，函数 $f(x)=x^2-2$ 在 $x=1$ 处为负、在 $x=2$ 处为正，但它\textbf{在有理数范围内永远不等于零}——“被经过的那个中间值”根本不存在！直觉在这里失灵，不是函数的问题，而是\textbf{数的世界本身有窟窿}。有理数像一件满是针眼的衬衫，看着完整，凑近了全是洞。

\begin{fanli}
警报一：函数 $f(x)=x^2-2$ 在有理数范围内“从负变正却不经过零”。所以介值定理的秘密一半在函数身上，另一半在实数身上——实数没有窟窿，这叫“完备性”。

警报二：“连续函数的图像一定光滑”也是错的。数学家魏尔斯特拉斯构造过一个函数，处处连续、处处没有切线——一条永远“毛茸茸”、任何一点都放大不出直线的曲线。它处处不抬笔，却处处扎手。可见“连续”比“光滑”弱得多，也微妙得多。
\end{fanli}

第三件事：就算函数不跳跃、数的世界也没有窟窿，直觉还可能败在“无穷”手里。函数 $g(x)=1/x$ 在区间 $(0,1]$ 上每一点都很乖，可是当 $x$ 靠近 $0$ 时，$g(x)$ 冲向无穷大，谁也拦不住。这条曲线在有限的范围里“装着”无限的高度。直觉说“一段有限长的曲线上总有个最高点”，在这里又错了。

三次被骗，教训是同一个：分析里的每一句话，都必须经得起最刁钻的抬杠。下面我们来一一修补。

\section{把“连续”说清楚}

先修第一个窟窿。到底什么叫连续？我们要的是这样一句话：\textbf{输入变化很小，输出变化就很小}。再准确一点：

\begin{dingyi}[连续（直观版）]
函数 $f$ 在点 $a$ 处\textbf{连续}，是指：当 $x$ 无限接近 $a$ 时，$f(x)$ 无限接近 $f(a)$。也就是说，你随便指定一个误差要求，比如“输出偏差不得超过 $0.001$”，我总能找到 $a$ 附近的一小段输入范围，保证这段范围内所有 $x$ 的输出都满足你的要求。函数在一个区间上连续，是指它在区间的每一点都连续。
\end{dingyi}

注意这个定义的“抬杠结构”：误差要求是\textbf{你}提的，可以多苛刻都行；而我\textbf{必须}给出应对方案。这就把“感觉上差不多”变成了可以逐条检验的协议。跳跃函数 $f(x)$ 在 $x=0$ 处过不了这一关：你要求“输出偏差小于 $1/2$”，可无论我把 $0$ 附近的输入圈缩得多小，圈里总有正数和负数，输出一会儿是 $1$ 一会儿是 $-1$，相差 $2$，永远超标。

\begin{shuyu}{连续函数}
\textbf{直观解释}：图像不抬笔就能画完；输入的小抖动只引起输出的小抖动。

\textbf{正式说法}：$x$ 趋向 $a$ 时，$f(x)$ 趋向 $f(a)$，即“极限值正好等于函数值”。

\textbf{一个例子}：$f(x)=x^2$ 处处连续，因为 $x$ 靠近 $a$ 时 $x^2$ 靠近 $a^2$；而跳跃函数在 $x=0$ 处不连续。
\end{shuyu}

第二个窟窿靠实数来补。实数系有一条宝贵性质：\textbf{任何“无限逼近某个位置”的过程，都能在实数里找到那个位置}。这叫做\textbf{完备性}。有理数不具备它——一列有理数 $1,\ 1.4,\ 1.41,\ 1.414,\ \dots$ 无限逼近 $\sqrt{2}$，可 $\sqrt{2}$ 不在有理数里。实数则没有这种漏洞。承认这一点之后，直觉的核心内容就可以被证明了。

\section{介值定理：不许跳，就得过}

\begin{dingli}[介值定理]
设 $f$ 在闭区间 $[a,b]$ 上连续，$f(a)<0$ 而 $f(b)>0$（或反过来）。那么区间里至少存在一点 $c$，使 $f(c)=0$。
\end{dingli}

更一般地，$f$ 会取到 $f(a)$ 与 $f(b)$ 之间的每一个值——海拔从 $100$ 米连续变到 $1500$ 米，必然恰好经过 $618$ 米，甚至恰好经过 $\sqrt{500}$ 米对应的高度。

证明的思路就是计算机里的“二分查找”：不断把嫌疑区间砍成两半，零点无处可逃。

\begin{proof}
设 $f(a)<0<f(b)$。第一步：看中点 $m_1=(a+b)/2$。若 $f(m_1)=0$，已经找到；若 $f(m_1)>0$，零点只能在左半边 $[a,m_1]$ 里；若 $f(m_1)<0$，零点只能在右半边 $[m_1,b]$ 里。无论哪种情况，我们都得到一个长度减半的新区间，它左端函数值为负、右端为正。

把这个操作重复下去：第 $n$ 步后，嫌疑区间的长度是 $(b-a)/2^n$，左端 $a_n$ 处 $f(a_n)\le 0$，右端 $b_n$ 处 $f(b_n)\ge 0$。左端点列 $a_1,a_2,a_3,\dots$ 单调不减、又被 $b$ 压着，右端点列单调不增、又被 $a$ 托着，两列点之间的距离趋向 $0$。由实数的完备性，两列点无限逼近同一个位置 $c$。

最后一步用连续性：$x$ 趋向 $c$ 时函数值趋向 $f(c)$。左端点列的所有函数值都不超过 $0$，所以它们的“逼近目标” $f(c)$ 不能超过 $0$；右端点列的所有函数值都不小于 $0$，所以 $f(c)$ 不能小于 $0$。两头的夹逼给出 $f(c)=0$。
\end{proof}

\begin{center}
\begin{tikzpicture}[scale=0.9]
  \draw[->] (-0.3,0) -- (7,0) node[right] {$x$};
  \draw[->] (0,-2.2) -- (0,2.6) node[above] {$y$};
  \draw[thick,domain=0.6:6.2,smooth] plot (\x, {0.6*\x-1.2+0.3*sin(2*\x r)});
  \draw[dashed] (0.6,0) -- (6.2,0);
  \fill (2.5,0) circle (1.5pt) node[below right] {$c$};
  \node[below] at (0.6,-0.65) {$a$};
  \node[above] at (6.2,2.5) {$b$};
  \node[left] at (0,2.2) {$f(b)>0$};
  \node[left] at (0,-1.6) {$f(a)<0$};
\end{tikzpicture}
\end{center}

这个证明值得多看一眼：它不仅证明零点\textbf{存在}，还真的\textbf{算}出了零点——每切一刀，误差减半，二十刀就把误差压到原来的百万分之一。用二分法求 $\sqrt{2}$：$1^2<2<2^2$，中点 $1.5$ 的平方是 $2.25>2$，往左切；中点 $1.25$ 的平方是 $1.5625<2$，往右切……十步之后你就有 $1.414$，和手算开方的结果吻合。分析的第一条定理，立刻变成了一台计算器。

\section{最高点与“脾气均匀”的函数}

介值定理管“中间值”，下面这条定理管“最值”，它们是一对兄弟。

\begin{dingli}[最大值定理]
闭区间 $[a,b]$ 上的连续函数，一定能取到最大值和最小值。也就是说，存在区间里的点 $p$ 和 $q$，使 $f(p)$ 不小于一切函数值，$f(q)$ 不大于一切函数值。
\end{dingli}

它的两个条件——“闭区间”和“连续”——一个都不能少。$g(x)=1/x$ 在 $(0,1]$ 上连续但区间不闭，它没有最大值；跳跃函数在闭区间上不连续，也可以把最高点“跳丢”。这条定理是所有“最优问题”的地基：你之所以敢问“怎样用料最省”“怎样利润最大”，是因为合理的函数在合理的范围里确实有答案，剩下的只是怎么找——那是第 10 章导数的活儿。

还有一个更精细的概念，考试不常考，却是分析的命脉：\textbf{一致连续}。普通连续是“每个点各自有一份应对方案”，点不同，方案可以不同；一致连续则要求\textbf{全区间共用一份方案}：你给定误差要求后，我能拿出一把统一长度的“尺子”，无论把它放在区间哪里，尺子两端之内的输入，输出偏差都达标。

为什么有的函数做不到？看 $g(x)=1/x$ 在 $(0,1]$ 上：越靠近 $0$，曲线越陡，同样的输入间隔 $0.001$，在 $x=1$ 附近输出只差大约 $0.001$，在 $x=0.001$ 附近输出能差上千。一把固定长度的尺子，越往左越不够用，永远找不到全区间通用的方案。而 $f(x)=x^2$ 在任何\textbf{闭}区间 $[a,b]$ 上都很老实：它的“坡度”有上限（不超过 $2$ 倍区间端点中较大的那个的绝对值），统一方案随手可得。可以证明：闭区间上的连续函数一定一致连续。直觉说“有限的舞台上演不出无限的脾气”，这次直觉赢了——但赢要靠证明，不靠感觉。

\section{柯西序列：不知道答案，也能证明答案存在}

介值定理的证明里藏着一个更深的问题。二分法造出了一列数 $a_1,a_2,\dots$，我们断言它“无限逼近某个位置 $c$”。可如果 $c$ 就是我们要找的东西，在找到它\textbf{之前}，凭什么知道它存在？难道不会逼近一个“空位置”吗，就像有理数逼近 $\sqrt{2}$ 那样？

十九世纪的数学家柯西给出了绝妙的回答：判断一列数“会收敛”，\textbf{不需要提前知道它收敛到谁}，只需要看这列数内部够不够团结。

\begin{dingyi}[柯西序列（直观版）]
一列数 $a_1,a_2,a_3,\dots$ 称为\textbf{柯西序列}，是指：你随便指定一个误差要求，这列数从某一项开始，\textbf{任意两项}之间的距离都不超过它。
\end{dingyi}

收敛的数列一定是柯西序列：大家都往同一个目标挤，彼此自然越靠越近。深刻的反问题是：柯西序列一定有目标吗？在有理数世界里不一定——$\sqrt{2}$ 的逼近列就是反例；在实数世界里，答案是\textbf{一定有}。这条“柯西序列必收敛”的性质，就是实数完备性的标准表述，也是整个分析大厦的地基石。

\begin{shuyu}{柯西序列}
\textbf{直观解释}：一列数“越到后来越抱团”，队尾的成员彼此挤得要多近有多近。

\textbf{正式说法}：任意给定误差 $\varepsilon>0$，存在一项，使得它后面的任意两项之差都小于 $\varepsilon$。

\textbf{一个例子}：$1,\ 1.4,\ 1.41,\ 1.414,\ \dots$（$\sqrt{2}$ 的逐位小数）是柯西序列：从第 $n$ 位小数之后，任意两项之差不超过 $10^{-n}$。它收敛的目标是 $\sqrt{2}$——这件事在实数里成立，在有理数里不成立。
\end{shuyu}

从此，“这个方程有解吗”有了新的答法：构造一列越来越团结的近似解，由完备性保证真正的解存在，再用连续性确认它真的满足方程。数学家用这套“先逼供、后定罪”的流程，证明了大量方程有解，哪怕解永远写不出精确表达式。

\section{幂级数：把函数拆成无限长的多项式}

如果说前面三节是“防守”——修补直觉的漏洞，那接下来的三节是“进攻”——用无穷过程造出惊人的新东西。第一件武器你已经见过雏形：无限小数。$\pi=3.14159\dots$ 其实是把 $\pi$ 拆成“$3$ 加上 $0.1$ 加上 $0.04$ 加上……”，每一位都是一个更小的修正项。幂级数干的是同一件事，只不过被拆的不是数，而是\textbf{函数}。

最基础的例子是等比数列求和。你知道 $1+x+x^2+\cdots+x^n=\dfrac{1-x^{n+1}}{1-x}$。当 $|x|<1$ 时，$x^{n+1}$ 趋向 $0$，于是
\[
\frac{1}{1-x}=1+x+x^2+x^3+\cdots \qquad(|x|<1).
\]
取 $x=1/2$ 验证：右边是 $1+1/2+1/4+\cdots=2$，左边是 $1/(1-1/2)=2$。吻合。

这个公式的读法是颠覆性的：函数 $\frac{1}{1-x}$ 居然等于一条“无限长的多项式”。为什么要这么拆？因为多项式是最好算的函数——只会乘法和加法。把一个“难算”的函数拆成无穷多个幂次的叠加，我们就有了它的“配方表”，想算多精确就取多少项。

最著名的配方是指数函数的：
\[
e^x=1+x+\frac{x^2}{2!}+\frac{x^3}{3!}+\frac{x^4}{4!}+\cdots
\]
它对\textbf{所有} $x$ 都成立（这和 $\frac{1}{1-x}$ 只在 $|x|<1$ 成立不同——每个配方表有自己的“适用半径”，称为收敛半径）。取 $x=1$，前六项给出
\[
1+1+\tfrac12+\tfrac16+\tfrac1{24}+\tfrac1{120}\approx 2.7167,
\]
而 $e\approx 2.71828$，五项就逼近了前三位小数。你计算器里那个 $e$ 的按键，内部多半就是这类级数在干活。

为什么适用半径有的大有的小？看 $\frac{1}{1-x}$ 就明白了：它在 $x=1$ 处“爆炸”（分母为零），所以它的配方表最远只能延伸到离爆炸点距离为 $1$ 的范围，再远就救不回来了。而 $e^x$ 在整个数轴上都风平浪静，配方表自然畅通无阻。粗略地说：\textbf{幂级数的适用范围，延伸到函数最近的“坏点”为止}。这条直觉在复数平面上会被证明得清清楚楚——“坏点”有时还藏在实轴看不见的地方，这也是数学家后来非研究复变函数不可的原因之一。

幂级数还有一个和导数有关的美妙性质：$e^x$ 的配方里，每一项求导后正好是前一项——$\frac{x^3}{3!}$ 求导得 $\frac{x^2}{2!}$，$\frac{x^2}{2!}$ 求导得 $x$，常数项 $1$ 求导得 $0$。整条级数求导后\textbf{原封不动还是自己}！这正是第 10 章说的“$e^x$ 的导数是它自己”在配方表上的倒影。

\section{傅里叶级数：任何声音都是一堆纯音}

第二件武器更浪漫。幂级数用“幂次”$1,x,x^2,\dots$ 做积木，法国数学家傅里叶提出了另一套积木：\textbf{正弦波}。他声称：几乎任何周期函数，都能用一串不同频率的正弦波叠加出来。这个主张当年被许多大数学家反对——用光滑的波浪拼出有棱角的图形，听起来像骗局。

请看这个著名的配方，一个方波（在 $0$ 和 $2\pi$ 之间，函数值在 $1$ 与 $-1$ 之间来回跳的矩形波）：
\[
\text{方波}(x)=\frac{4}{\pi}\left(\sin x+\frac{\sin 3x}{3}+\frac{\sin 5x}{5}+\frac{\sin 7x}{7}+\cdots\right).
\]
只取第一项，是一条起伏平缓的正弦曲线，和方波相去甚远；加上第二项，波峰开始变平；加到第五项，矩形的样子已经很像了，只剩边缘的细微波纹。取的项越多，波纹越小，棱角越锋利。

这个配方里藏着一条朴素的规律：每个正弦波的频率是前一个的“下一个奇数倍”，而振幅按频率的倒数衰减——三倍频率配三分之一高度，五倍频率配五分之一高度。为什么高频的成分越来越少？因为方波的大部分“体形”是平缓的，只有棱角的细节才需要高频波来雕刻；细节总是比主体少。这条规律在所有傅里叶分解里都成立：\textbf{低频决定轮廓，高频决定细节}。记住这句话，待会儿的实验和后面的连接线里都要用到。

傅里叶当年把这套想法写成论文提交给巴黎的科学院，评审的数学家们（包括以严格著称的拉格朗日）拒绝签字：他们要求先回答“这串无穷叠加到底收不收敛、收敛到什么”，而这正是本章前半部分的柯西序列和完备性要回答的问题。历史在这里绕了一个圈：为了审判傅里叶的级数，数学家们被迫把“连续”“收敛”这些词说得更清楚——分析的严格化，有一半是被这些波浪逼出来的。

\begin{center}
\begin{tikzpicture}[scale=0.9]
  \draw[->] (-0.3,0) -- (6.8,0) node[right] {$x$};
  \draw[->] (0,-1.9) -- (0,1.9);
  \draw[gray,dashed] (0,1.25) -- (3.14,1.25) (3.14,-1.25) -- (6.28,-1.25);
  \draw[gray,dashed] (3.14,1.25) -- (3.14,-1.25) (6.28,1.25) -- (6.28,-1.25);
  \draw[domain=0:6.28,samples=80,smooth] plot (\x, {1.273*sin(\x r)});
  \draw[thick,domain=0:6.28,samples=120,smooth] plot (\x, {1.273*(sin(\x r)+sin(3*\x r)/3+sin(5*\x r)/5)});
  \node[right] at (6.3,1.35) {五项叠加};
  \node[right] at (6.3,0.75) {只有 $\sin x$};
\end{tikzpicture}
\end{center}

代一个具体的数进去，你会得到一个百年不衰的惊喜。取 $x=\pi/2$：方波在这一点取值 $1$；而 $\sin(\pi/2)=1$，$\sin(3\pi/2)=-1$，$\sin(5\pi/2)=1$，交替正负。于是
\[
1=\frac{4}{\pi}\left(1-\frac13+\frac15-\frac17+\cdots\right),
\quad\text{即}\quad
\frac{\pi}{4}=1-\frac13+\frac15-\frac17+\cdots .
\]
圆的周长里藏着的 $\pi$，竟然能用全体奇数的倒数交错相加算出来。算到第 $100$ 项，你大约能得到 $\pi$ 的两位小数——收敛虽慢，但谁能否认这个等式的漂亮？

\begin{shiyan}
\textbf{亲手拼一个方波。}在方格纸上，先画 $y=\sin x$（$x$ 从 $0$ 到 $2\pi$，取 $x=0, 0.5, 1, \dots$ 用计算器算值描点）。再画 $\frac{1}{3}\sin 3x$——频率是三倍、高度是三分之一。把两条曲线在相同 $x$ 处的高度\textbf{相加}，描出新曲线。你会看到波峰被压平、波谷被抬起，矩形的雏形出现了。有余力再加 $\frac{1}{5}\sin 5x$。三种颜色叠在一张图上，就是傅里叶两百年前看到的景象。

\textbf{另一个实验}：对着手机上的频谱软件（或音乐软件的“均衡器”界面）哼一个长音。屏幕上会亮起几根柱子：最低最亮的是基音，上面是泛音。你的嗓音和钢琴弹同一个音，听起来完全不同，差别就在泛音柱子的相对高度上——声音里的“音色”，就是傅里叶系数。
\end{shiyan}

\begin{lianjie}
傅里叶级数就在你口袋里。音乐文件（如 MP3）的压缩原理，是把声音切成小段，每段分解成不同频率正弦波的叠加，再把人耳听不见的高频分量删掉——数据量大减，听感几乎不变。图片压缩（JPEG）用的是同一家族的“余弦变换”：一张照片中平缓变化的区域，低频系数就能描述，细节才需要高频。电话降噪、地震波分析、医学成像，全是傅里叶思想的子孙。一个当年被嘲笑“不可能”的主张，如今每天处理着全人类的图片和声音。
\end{lianjie}

\section{复数平面与欧拉公式}

第三件武器，来自前面章节认识的老朋友：虚数单位 $i$，满足 $i^2=-1$。当时你可能觉得它是个工具人——为了让所有二次方程有解而请来的临时工。这一节你会看到，$i$ 一旦和分析联手，会产生数学里公认最美的一个公式。

先把复数摆在平面上：$a+bi$ 对应横坐标 $a$、纵坐标 $b$ 的点。加法就是向量的平移，这很自然。妙的是乘法：乘以 $i$，相当于把向量\textbf{逆时针旋转九十度}。因为 $1$ 乘 $i$ 得 $i$——点 $(1,0)$ 转到了 $(0,1)$；再乘 $i$ 得 $-1$——又转九十度。转两次共一百八十度，正是“乘以 $-1$”。所谓 $i^2=-1$，用几何语言说就是“两个九十度等于一百八十度”。复数乘法原来是\textbf{旋转加伸缩}。

现在请幂级数登场。数学家欧拉做了一件看似疯狂的事：把纯虚数 $i\theta$ 代入 $e^x$ 的配方表，看会发生什么。利用 $i^2=-1$、$i^3=-i$、$i^4=1$ 的循环，奇数项和偶数项自动分成两队：
\[
e^{i\theta}=\left(1-\frac{\theta^2}{2!}+\frac{\theta^4}{4!}-\cdots\right)
+i\left(\theta-\frac{\theta^3}{3!}+\frac{\theta^5}{5!}-\cdots\right).
\]
两个括号里的级数，恰好就是 $\cos\theta$ 和 $\sin\theta$ 的配方表（可以用第 10 章的导数知识验证：它们的各阶导数在 $0$ 点的值完全吻合）。于是：

\begin{dingli}[欧拉公式]
对任意实数 $\theta$，
\[
e^{i\theta}=\cos\theta+i\sin\theta .
\]
\end{dingli}

它的几何意义令人屏息：右边是单位圆上角度为 $\theta$ 的点，左边是“$e$ 的纯虚数次幂”。也就是说，$e^{i\theta}$ 随着 $\theta$ 增大，在单位圆上匀速转圈。代入 $\theta=\pi$：$\cos\pi=-1$，$\sin\pi=0$，于是
\[
e^{i\pi}+1=0 .
\]
一个公式里站着五个最重要的常数：$0$、$1$、$\pi$、$e$、$i$，被加、乘、乘方三种运算缝在一起。数学家们投票选“最美公式”时，它几乎从不落选。

\begin{center}
\begin{tikzpicture}[scale=1.6]
  \draw[->] (-1.6,0) -- (1.8,0) node[right] {实轴};
  \draw[->] (0,-1.6) -- (0,1.8) node[above] {虚轴};
  \draw (0,0) circle (1);
  \draw[->,thick] (0,0) -- (0.5,0.866);
  \fill (0.5,0.866) circle (1pt) node[above right] {$e^{i\theta}$};
  \draw (0.28,0) arc (0:60:0.28) node[midway,right] {$\theta$};
  \fill (-1,0) circle (1pt) node[below left] {$-1=e^{i\pi}$};
  \fill (1,0) circle (1pt) node[below right] {$1$};
  \fill (0,1) circle (1pt) node[above left] {$i$};
\end{tikzpicture}
\end{center}

欧拉公式还是傅里叶级数的“官方语言”。用 $\cos\theta=\frac{e^{i\theta}+e^{-i\theta}}{2}$ 和 $\sin\theta=\frac{e^{i\theta}-e^{-i\theta}}{2i}$，所有正弦余弦的叠加都可以改写成 $e^{ik\theta}$ 的叠加，三角恒等式变成了简单的指数运算。现代工程师做傅里叶分析时，笔下几乎没有 $\sin$ 和 $\cos$，全是旋转的小箭头 $e^{i\theta}$。

再往前一步，就是\textbf{复分析}：把 $x$ 换成复数 $z$，研究复变函数 $f(z)$。这里发生了数学中少见的奇迹：在复平面上，“可导”这个条件变得极其强大——一个函数只要在某个区域里处处有导数，它就自动光滑、自动能展开成幂级数，函数值还被边界上的取值完全锁死。复变函数像一块烧红的铁板，动一处而全身知。这种刚性让复分析既优雅又有力，也让它长出了数学里最深的问题之一。

\section{本章的收获}

\begin{toolbox}
本章新增的工具：

\begin{itemize}
  \item \textbf{连续（直观版）}：你提误差要求，我给应对方案——输入足够接近，输出就足够接近。
  \item \textbf{实数完备性}：柯西序列（越来越抱团的数列）在实数里必有极限。有理数没有这条性质，这是两套数的根本差别。
  \item \textbf{介值定理}：连续函数取遍两端之间的一切值；证明方法是二分法，顺便还能算零点的近似值。
  \item \textbf{最大值定理}：闭区间上的连续函数必有最大值和最小值；条件“闭区间”“连续”缺一不可。
  \item \textbf{一致连续}：全区间共用一份应对方案；闭区间上的连续函数一定一致连续。
  \item \textbf{幂级数}：把函数拆成无限长的多项式，如 $e^x=\sum x^n/n!$；每个配方有适用半径。
  \item \textbf{傅里叶级数}：把周期函数拆成正弦波的叠加；方波的展开顺带算出 $\pi/4=1-1/3+1/5-\cdots$。
  \item \textbf{欧拉公式}：$e^{i\theta}=\cos\theta+i\sin\theta$，复数乘法是旋转加伸缩，特例 $e^{i\pi}+1=0$。
\end{itemize}
\end{toolbox}

\begin{pandeng}
继续攀登的路标：

\begin{itemize}
  \item \textbf{$\varepsilon$--$\delta$ 语言}：本章“你提要求、我给方案”的协议，在大学里会写成精确的符号形式，这是数学分析的第一课。
  \item \textbf{魏尔斯特拉斯函数}：处处连续、处处不可导的“毛刺曲线”，它的严格构造用到了一致收敛的函数项级数。
  \item \textbf{泰勒级数}：幂级数配方的一般理论——函数的各阶导数决定它的全部系数。
  \item \textbf{快速傅里叶变换（FFT）}：把傅里叶分解的计算量降到不可思议程度的算法，被称为二十世纪最重要的算法之一。
  \item \textbf{黎曼 $\zeta$ 函数与黎曼猜想}：复分析皇冠上的明珠，它把本章的级数和下一章的质数缝在一起——悬赏一百万美元的未解之谜。
\end{itemize}
\end{pandeng}

\section*{思考题}

\textbf{观察题}
\begin{sikao}
  \item 某天中午 $12$ 点你从山脚出发上山，下午 $4$ 点到顶；第二天中午 $12$ 点你从山顶沿同一条路下山，下午 $4$ 点到山脚。证明（用介值定理）：两天中必有同一时刻，你处在同一海拔。\emph{提示}：把两天的“海拔关于时刻”的函数画在同一坐标系里，考虑它们的差。
  \item 方程 $x^3+x-1=0$ 在 $0$ 与 $1$ 之间有解吗？能否用二分法手算三步，把解的位置缩小到长度 $1/8$ 的区间里？\emph{提示}：算 $x=0,\frac12,1$ 处的函数值，再切。
\end{sikao}

\textbf{证明题}
\begin{sikao}
  \setcounter{sikaoi}{2}
  \item 证明：连续函数 $f$ 若满足 $f(a)\cdot f(b)<0$，则它在 $(a,b)$ 内必有零点——这正是介值定理。请把正文二分法证明中的“左端点列单调不减且有上界，故逼近某点 $c$”这一步，用自己的话补全理由。\emph{提示}：一列数越排越靠右、却永远越不过一堵墙，它会怎么样？
  \item 设数列 $a_n=\frac{n}{n+1}$（即 $\frac12,\frac23,\frac34,\dots$）。验证它是柯西序列：当 $m,n$ 都大于某个 $N$ 时，$|a_m-a_n|$ 能有多小？\emph{提示}：$|a_m-a_n|=\left|\frac{1}{n+1}-\frac{1}{m+1}\right|$，分别给两个分数找上界。
\end{sikao}

\textbf{挑战题}
\begin{sikao}
  \setcounter{sikaoi}{4}
  \item 用 $e^{i\theta}=\cos\theta+i\sin\theta$ 计算 $e^{i\alpha}\cdot e^{i\beta}$，两边分别用指数律和复数乘法展开，比较实部与虚部，你能“免费”推出哪两个三角恒等式？\emph{提示}：$(\cos\alpha+i\sin\alpha)(\cos\beta+i\sin\beta)$ 展开后，实部是什么？
  \item 级数 $1-\frac13+\frac15-\frac17+\cdots$ 收敛到 $\pi/4$，但它收敛得慢：前 $100$ 项的误差大约有多大？观察相邻两项的包夹关系来估计。\emph{提示}：交错级数的部分和在目标值两侧来回跳，误差不超过被丢掉的头一项。
\end{sikao}

\section*{通往下一章的门}

这一章的结尾，留给级数一个意味深长的镜头。考虑调和级数
\[
1+\frac12+\frac13+\frac14+\frac15+\cdots .
\]
它的每一项都趋向 $0$，可它的和却\textbf{趋向无穷大}——把项分组：$\frac13+\frac14>\frac12$，$\frac15+\frac16+\frac17+\frac18>\frac12$，$\frac19$ 到 $\frac1{16}$ 的八项之和又大于 $\frac12$……你能凑出无穷多个“半个”，和当然拦不住。

欧拉盯上了这件小事，并追问了一个改变数学史的问题：调和级数的发散，和\textbf{质数}有什么关系？他发现了一个惊人的恒等式：全体正整数的倒数之和，竟然等于只由质数造出来的一个无穷乘积。如果质数只有有限个，这个乘积就是个有限的数，可调和级数明明是发散的——矛盾！于是“质数有无穷多个”这件两千年前欧几里得就证明过的事，被欧拉用分析的语言又证明了一遍，而且锋利得多。

这扇门的背后是一个新世界：整数的乘除机关里，藏着级数、函数甚至复平面的影子。而想读懂这一切，你得先回到最朴素的问题——质数到底是什么？它们按什么规律排列？为什么 $7\times 11\times 13=1001$ 这样的分解是唯一的？下一章，我们放下连续与无限，走进整数的世界：质数、同余，和它们守护的现代密码。
